\documentclass[a4paper,11pt]{article}
\usepackage{exam-style-842-unit}
\examinehead{842 信息化概论 · 第 05 单元练习册}{练习一 参考答案与解析}
\begin{document}

\unitanswerscover{842 信息化概论 · 单元练习册}{第 05 单元参考答案}{练习一（云计算基础）的答案与解析；每道判断题均给出改正后的正确表述，教材未展开的内容均注明来源。}

\answerpart{一、单项选择题（10 题，每题 3 分，共 30 分）}

\q A。\begin{analysis}早在 20 世纪 60 年代，麦卡锡（John McCarthy）就提出了把计算能力作为一种像水和电一样的公用事业提供给用户（教材 4.1.1，第 43 页）。B 施密特是 2006 年首次提出“云计算”概念的人；C、D 与本题表述无关。\end{analysis}

\q A。\begin{analysis}2006 年 8 月 9 日，Google 首席执行官埃里克·施密特（Eric Schmidt）在搜索引擎大会（SES San Jose 2006）上首次提出“云计算”（Cloud Computing）的概念（教材 4.1.1，第 43 页）。B 贝索斯是亚马逊创始人，2011 年前后提出“弹性计算云”，教材只述亚马逊公司行为；C 卡廷是 Hadoop 项目创建者；D 扎克伯格与 Facebook 相关。\end{analysis}

\q A。\begin{analysis}NIST 的定义是：云计算是一种按使用量付费的模式，这种模式提供可用的、便捷的、按需的网络访问，进入可配置的计算资源共享池（资源包括网络、服务器、存储、应用软件、服务），这些资源能够被快速提供（教材 4.1.2，第 44 页）。B、C、D 都不是计费依据。\end{analysis}

\q B。\begin{analysis}云计算支持用户在任意位置、使用各种终端获取应用服务；用户请求的资源来自“云”，而不是固定的有形的实体，用户无需了解、也不用担心应用运行的具体位置（教材 4.1.3“虚拟化”条第 44 页）。A 指“云”的服务器规模大；C 指多副本容错、计算节点同构可互换；D 指不针对特定应用、同一朵云可支撑不同应用。\end{analysis}

\q C。\begin{analysis}“云”的规模可以动态伸缩，满足应用和用户规模增长的需要，这是高可扩展性（教材 4.1.3 第 44 页）。A 通用性指云计算不针对特定应用；B 高可靠性靠数据多副本容错等措施；D 廉价指用廉价节点构成云、自动化集中管理降低成本。\end{analysis}

\q D。\begin{analysis}NIST 把云计算的基本特征概括为按需自助服务、广泛网络接入、资源池化、快速弹性、可计量服务 5 个（NIST 定义，属通行标准表述，教材 4.1.2 只引用了 NIST 的定义而未单列这 5 条；教材 4.1.3 自己归纳的是 7 个特点）。D 数据独占性与“资源共享池”相矛盾，云计算恰恰强调多用户共享资源池。\end{analysis}

\q B。\begin{analysis}私有云系统存在于企业防火墙之内，只为企业内部服务；与公有云相比安全性更好，但成本更高，规模经济效益受到限制（教材 4.1.4，第 46 页）。A 公有云是服务商为公众提供的平台；C 混合云同时提供公有和私有服务；D 社区云为特定社区提供，教材未列。\end{analysis}

\q C。\begin{analysis}按 NIST 的权威定义，云计算的服务模式有 IaaS（基础设施即服务）、PaaS（平台即服务）和 SaaS（软件即服务）三大类或层次，由低到高依次排列（教材 4.1.4 如图 4.1，第 45—46 页）。A、B、D 都颠倒了层次顺序。\end{analysis}

\q D。\begin{analysis}教材列出公有云的 4 个致命缺陷：（1）政府没有运营云的法律法规，数据出了问题无法追究责任；（2）云运营商计费标准不统一，流量计费千差万别；（3）能给集团企业提供云服务的运营商数量还不多；（4）一旦更换云服务商，数据迁移是个大问题（教材 4.1.4，第 45 页）。D 的表述与第（4）条相反，且私有云迁移难度并不比公有云小。\end{analysis}

\q A。\begin{analysis}教材 4.1.5 指出：虚拟化是云计算最重要的核心技术之一，为云计算服务提供基础架构层面的支撑，是 ICT 服务快速走向云计算的最主要驱动力（第 47 页）。B 绿色节能技术、D 负载均衡都是相关技术，但题目问的是“基础架构层面支撑”的最重要核心技术；C 分布式存储属于大规模数据处理技术的内容。\end{analysis}

\answerpart{二、多项选择题（4 题，每题 4 分，共 16 分）}

\q ABC。\begin{analysis}教材 4.1.3 列出的云计算特点为：超大规模、虚拟化、高可靠性、通用性、高可扩展性、按需服务、廉价（第 44 页）。A、B、C 均在列；D 数据独占性错误，云是共享资源池，资源按需分配而非独占。\end{analysis}

\q ACD。\begin{analysis}A 正确：公有云理论上是服务商为公众提供、任何人授权即可接入的平台（教材 4.1.4，第 45 页）。C 正确：私有云安全性更好，但成本更高，基础设施利用率远低于公有云（第 46 页）。D 正确：混合云是同时提供公有和私有服务的折中方案（第 46 页）。B 错误：私有云存在于企业防火墙之内，只为企业内部服务，不是“防火墙之外、多个外部客户共享”。\end{analysis}

\q ABC。\begin{analysis}教材 4.1.5 指出云计算的核心技术“以虚拟化技术、大规模数据处理技术、云计算平台管理技术、绿色节能技术最为关键”（第 47—48 页）。A、B、C 分别对应前三项；D 磁盘碎片整理是单机操作系统的磁盘维护功能，不属于云计算核心技术。\end{analysis}

\q ABC。\begin{analysis}教材 4.1.5 指出：云计算是网格计算、分布式计算、并行计算、效用计算、网络存储、虚拟化、负载均衡等传统计算机技术和网络技术发展融合的产物（第 47 页）。A、B、C 均在列；D 生物计算不属于教材所列的云计算来源技术（生物计算与生物信息学相关，不是云计算的技术基础）。\end{analysis}

\answerpart{三、判断题（6 题，每题 2 分，共 12 分）}

\q √。\begin{analysis}2006 年 8 月 9 日，Google 首席执行官埃里克·施密特在搜索引擎大会（SES San Jose 2006）上首次提出“云计算”概念（教材 4.1.1，第 43 页）。\end{analysis}

\q ×。改正：虚拟化是云计算的重要组成部分，但不是全部；云计算还包括大规模数据处理技术、云计算平台管理技术、绿色节能技术等核心技术。\begin{analysis}教材 4.1.5 明确批评了“云计算就是虚拟化”的误区：“虚拟化是云计算的重要组成部分但不是全部”（第 47 页）。\end{analysis}

\q √。\begin{analysis}教材 4.1.4 从云的建设地点、云服务的协议开发程度、服务对象 3 个方面区分公有云与私有云（第 46 页）。\end{analysis}

\q ×。改正：按 NIST 的权威定义，云计算的服务模式分为 IaaS、PaaS、SaaS 三大类（层次），没有 DaaS。\begin{analysis}教材 4.1.4 指出“云计算的服务模式有 IaaS、PaaS 和 SaaS 这 3 个大类或层次，这是目前被业界最广泛认同的划分”（第 45 页）。\end{analysis}

\q √。\begin{analysis}私有云系统存在于企业防火墙之内，只为企业内部服务；与公有云相比安全性更好，但成本更高，规模经济效益也受到限制（教材 4.1.4，第 46 页）。\end{analysis}

\q ×。改正：云计算不是一种全新共享基础架构的方法，而是网格计算、分布式计算、并行计算、效用计算、网络存储、虚拟化、负载均衡等传统计算机技术和网络技术发展融合的产物。\begin{analysis}教材 4.1.5 开篇即指出“云计算不是一种全新共享基础架构的方法”，并列举其来源技术（第 47 页）。\end{analysis}

\answerpart{四、配对题（6 对，每对 1 分，共 6 分）}

\q （1）—B\quad （2）—A\quad （3）—C\quad （4）—D\quad （5）—F\quad （6）—E。
\begin{analysis}麦卡锡在 20 世纪 60 年代提出计算能力应像水和电一样作为公用事业（B）；施密特 2006 年首次提出“云计算”概念（A）；1999 年 Salesforce.com 提出通过网站向企业提供企业级应用（C）；亚马逊 2002 年提供包含存储、计算能力等资源的 Web Service，2005 年提出弹性计算云 EC2（D）；IaaS 位于最底层，提供服务器、存储、网络等基础设施服务（F）；SaaS 位于第三层，用户通过网页浏览器使用远程服务器上的应用（E）。（教材 4.1.1、4.1.4，第 43、45—46 页）\end{analysis}

\answerpart{五、简答题（3 题，每题 8 分，共 24 分）}

\q 参考答案要点：
（1）教材定义（2 分）：云计算是基于互联网的相关服务的增加、使用和交付模式，通常涉及通过互联网来提供动态易扩展且经常是虚拟化的资源；云是网络、互联网的一种比喻说法（教材 4.1.2，第 43—44 页）。
（2）NIST 定义（2 分）：云计算是一种按使用量付费的模式，这种模式提供可用的、便捷的、按需的网络访问，进入可配置的计算资源共享池（资源包括网络、服务器、存储、应用软件、服务），这些资源能够被快速提供，只需投入很少的管理工作，或与服务供应商进行很少的交互。
（3）七个特点（4 分，答对 5 点即可得满分）：超大规模——Google 云已有 100 多万台服务器，企业私有云一般也有数百上千台；虚拟化——用户可在任意位置、用各种终端获取服务，无需了解资源具体位置；高可靠性——数据多副本容错、计算节点同构可互换；通用性——不针对特定应用，同一朵云可同时支撑不同应用；高可扩展性——规模可动态伸缩；按需服务——资源池按需购买，像水电煤气一样计费；廉价——可用廉价节点构成云，自动化集中管理降低数据中心成本（教材 4.1.3，第 44 页）。
（4）答完定义后若能补充 NIST 概括的按需自助服务、广泛网络接入、资源池化、快速弹性、可计量服务 5 个基本特征，可在本小题内作为加分说明（该 5 条为 NIST 通行表述，教材未单列）。

\q 参考答案要点（教材 4.1.5，第 47—48 页）：
（1）虚拟化技术（2 分）：为云计算服务提供基础架构层面的支撑，是 ICT 服务快速走向云计算的最主要驱动力；它是在软件中仿真计算机硬件、以虚拟资源为用户提供服务的计算形式，旨在合理调配计算机资源、提高效率；其最大好处是增强系统的弹性和灵活性，降低成本、改进服务、提高资源利用效率。虚拟化是云计算的重要组成部分但不是全部。
（2）大规模数据处理技术（2 分）：云计算集中大量服务器产生规模效应，提供海量数据存储和处理能力；通常采用分布式存储技术把数据存放在不同物理设备中，可靠性高、扩展性好；从编程模式看，云计算是多用户、多任务、支持并发处理的系统，广泛采用分布式并行编程模式。
（3）云计算平台管理技术（2 分）：云采用分布式存储，必须引入分布式资源管理技术，保证多节点状态同步、单节点故障时其他节点不受影响；面对数百台至上万台服务器、跨地域的资源和数以千计的应用，需要统一管理、统一分配、统一部署、统一监控和统一备份，打破应用对资源的独占。
（4）绿色节能技术（2 分）：云计算以低成本、高效率著称，规模经济效益在提高资源利用率的同时节省大量能源；业界通过把机房建在北极圈旁、利用海水或湖水冷却、太阳能供电等方式降低数据中心能耗与 PUE，例如阿里巴巴新一代绿色数据中心设计年均 PUE 低于 1.3。

\q 参考答案要点：
（1）联系（3 分）：教材 4.1.5 指出，云计算不是全新的共享基础架构方法，而是网格计算、分布式计算、并行计算、效用计算、网络存储、虚拟化、负载均衡等传统计算机技术和网络技术发展融合的产物（第 47 页）；三者都利用网络把多台计算机组织起来协同工作，都以资源共享、提高利用率为目标。
（2）区别（5 分，教材未给网格计算、分布式计算下定义，以下按通行表述整理）：分布式计算强调把计算任务分解为若干子任务、分配到多台计算机上并行处理再汇总结果，重在“任务分解与并行”；网格计算强调把地理上分布、异构的计算资源集成为一台“虚拟超级计算机”，跨机构共享空闲资源、统一调度，多用于大规模科学计算；云计算强调以互联网为交付方式，把计算资源池化后按需、弹性地提供给用户，并按使用量付费，重在“服务化、弹性与按需自助”，用户无需关心资源位置和运维。
（3）结论（1 分）：三者不是替代关系而是继承与融合关系；云计算吸收了分布式计算的任务分解、网格计算的资源共享思想，并借助虚拟化、平台管理等技术把计算变成像水电一样可计量的服务。

\answerpart{六、综合题（1 题，共 12 分）}

\q 参考答案要点：
（1）部署模式（4 分）：院内涉密数据部分应采用私有云——私有云部署在企业防火墙内、只为企业内部服务，对数据、安全性和服务质量的控制最有效（教材 4.1.4，第 46 页）；面向社会公众的农技问答与行情服务应采用公有云——理论上任何人都可授权接入，能发挥规模经济效益（第 45 页）；两者组合即混合云，是同时提供公有和私有服务的折中方案，兼顾涉密数据安全与公众服务的弹性低成本（第 46 页）。
（2）服务模式（4 分）：底层计算资源（服务器、存储、网络）对应 IaaS，用户购买服务器基础服务即可，无需自行购买、托管和维护硬件（4.1.4，第 45 页）；开发运行环境对应 PaaS，平台提供分布式操作系统、分布式数据库等基础软件，各课题组的开发在这一层进行（第 45 页）；面向公众的问答服务对应 SaaS，用户通过网页浏览器即可使用远程服务器上的应用（第 45—46 页）。
（3）辨析（4 分）：该说法不准确。虚拟化确实是云计算最重要的核心技术之一，为云计算提供基础架构层面的支撑，但“虚拟化是云计算的重要组成部分但不是全部”（4.1.5，第 47 页），云计算还包括大规模数据处理技术、云计算平台管理技术、绿色节能技术。从来源看，云计算是网格计算、分布式计算、并行计算、效用计算、网络存储、虚拟化、负载均衡等技术融合的产物，并非买了虚拟化软件就完成了云计算的资源池化、按需自助服务、弹性伸缩与计量计费。与网格计算相比，云计算面向公众按需交付并按使用量收费，而网格计算多用于跨机构的科研资源共享；与分布式计算相比，云计算还要解决多租户隔离、平台管理、服务交付等问题。因此，虚拟化是必要的基础，但不是云计算的充分条件。

\end{document}
